% LSPV-4 REPORT::ARC-COVER-GATE-SCOPE262::2026-10-30T20:36:00-07:00
\documentclass[11pt]{article}
\usepackage[margin=.85in]{geometry}
\usepackage{amsmath,amssymb,amsthm,booktabs,hyperref,xurl}
\newtheorem{theorem}{Theorem}
\title{ARC-COVER-GATE-SCOPE262: Single-Cube Semantic Cover and the Actual FLMV Model}
\author{LSPV-4 Circuit Resource Auditor}
\date{30 October 2026, 20:36 PDT}
\begin{document}\maketitle
\noindent\textbf{Status.} PROVED/EXACT explicit restricted many-one reduction, mathematical peer review pending. KNOWN THEOREM coNP-completeness of entailment/backbone recognition and FLMV. COMPUTATIONAL matched finite controls. CONDITIONAL native model transfer. OPEN arbitrary-$A\mapsto\Phi(A)$ contraction and $P$ vs NP.
\section{One exact falsifiable gate}
The source-relative completeness requirement in the LSPV-5 ARC-RESIDUAL-COVER-GATE231 report is $F\models\bigvee_i\operatorname{Cube}(\rho_i)$ if selected exact residual answers are to support an UNSAT conclusion without a separate exact fallback. Falsifiers: a singleton cube always has a cheap semantic cover check (FALSE); one SAT child proves its cube covers the source models (FALSE); FLMV Theorem 3.4 literally states a completely arbitrary virtual-input-accessor simulation theorem (FALSE as a citation claim, not a counterexample to a separately proved simulation); a direct word-RAM bit accessor suffices for transfer to native FLMV space bounds (UNPROVED).
\begin{theorem}
Given a proper signed 3CNF $F$ and a distinguished variable $y$, with the promise that $F$ is satisfiable, deciding $F\models y$ is coNP-complete under polynomial-time many-one reductions.
\end{theorem}
\begin{proof}
For any proper 3CNF $H=\bigwedge_{j=1}^m(a_j\vee b_j\vee c_j)$ with $n$ source variables, introduce fresh $y,z_1,\dots,z_m$ and define
\[
F_H=\bigwedge_{j=1}^m[(y\vee a_j\vee z_j)\land(\neg z_j\vee b_j\vee c_j)].
\]
For each source assignment $x$, the $j$th pair under $y=0$ permits at least one value of $z_j$ precisely when $a_j\vee b_j\vee c_j$ is true. Thus
\[
(\exists z)\,F_H(x,0,z)\ \Longleftrightarrow\ H(x).
\]
Under $y=1$, every $z_j=0$ works regardless of $x$, so every $F_H$ is SAT. Therefore $F_H\models y$ iff $H$ is UNSAT. The mapping produces $2m$ proper 3-clauses and $n+m+1$ variables and has logarithmic-index streaming syntax. A falsifying witness for $F\models y$ is a complete assignment satisfying $F$ with $y=0$, so the promise problem is in coNP. This proves coNP-hardness and membership. This is a known SAT-backbone/entailment phenomenon, not a new general complexity lower bound.
\end{proof}
\section{Both-label ARC interpretation and charged barriers}
The identical source-to-output program handles SAT and UNSAT $H$. Its output is always SAT. For an UNSAT $H$, the sole cube $y=1$ is a valid source-relative model cover. For a SAT $H$, that cube misses genuine models with $y=0$, although its restricted child is trivially SAT. Thus complete child-answer correctness is distinct from a checked cover certificate. A solver may stop once a SAT witness is verified, or verify a genuine proof of coverage, or retain an exact uncovered-region fallback. This construction does not force any universal SAT solver to perform the semantic check.
A naive cover verifier calling an arbitrary correct SAT guest $A$ on $F_H\land\neg y$ at the parent input scale pays parent-size guest time. DELAY wrappers $A_D$ that spend $N^D$ steps before calling a correct decider preserve both labels but violate a purported bound $K'\le\max\{11/10,3/5+K/2\}$ on this restricted implementation for sufficiently large $D$. This is the previously known parent-call firewall applied to the cover gate, not a general lower bound on all adaptive policies.
The $2m$-clause reduction requires $\Theta(m\log(n+m))$ explicit output bits in a standard named-variable encoding. Physically buffering that output on an FLMV work tape incurs a linear-size counted workspace contribution. A virtual child requires all input-head, index-erasure, far-blank, ordinary-move, multi-tape, source-read and ancestor costs paid. A short local bit emitter on a word RAM does not by itself prove the same-machine input transfer. Any resulting native-FLMV contradiction is CONDITIONAL on an exact model-preserving transfer.
\section{Independent FLMV audit and citation scope}
For $f_0=1$, $f_k=cf_{k-1}/(1+f_{k-1})$,
\[
d_*(c)=\max_{k\ge1}\min\{(1/f_k-c)/k,c^{-k}\}.
\]
At $c=6/5$, $(d_1,d_2,d_3)=(7/15,23/45,401/810)$; all $k\ge4$ have $d_k\le625/1296<23/45$. At $c=241/200$, $(d_1,d_2,d_3)=(21919/48200,11642479/23232400,4077837439/8398512600)$; all $k\ge4$ have $d_k\le1600000000/3373402561<d_2$. Thus both exact maxima are attained at $k=2$, with the entire tail certified. The proposed formal fixed point $(6/5,1/2)$ would meet strict margins $1/200$ in time exponent and $26279/23232400$ in space exponent only under the still-open all-input hereditary transformation premises.
The ORIGINAL FLMV Theorem 3.4 states a nondeterministic-linear-time multitape to SAT simulation with a specific Cook-style locally generated SAT formula. It does not \emph{state} the fully generic guest/total-external-accessor theorem attributed to it in part of root main.tex. The underlying project head-shadow result may be valid as an independently proved lemma, under its complete model and accessor premises. LSPV-5 should correct attribution; no existing independent result is refuted here.
\section{Reproducible finite controls and audit limits}
A separate Python producer and independently implemented checker passed 135,508 comparison groups: 65,656 coherent original assignments across ten sources, 2,066 four-mode clause-pair checks, and all original transformed-clause S/U targets, inversions, actual global endpoints, multiple flips and source-parent prohibited cubes. W1/S15 are a same-support 15-variable/20-clause strict signed exact-(2,2) pair. W1 is UNSAT with zero models, S15 SAT with 861 models. Both transformed $F_H$ are SAT. The singleton-cover decisions respectively are YES and NO. Baseline W1 DPLL had 151 nodes and 11,555 selected operations; transformed $F_H\land\neg y$ had 151 nodes and 38,190 operations. S15 baseline had 9 nodes and 784 operations, while transformed checking had 12 nodes and 3,630 operations. XOR3, hitting, Horn, padded-Krom and two-block affine controls also passed. These are instrumented Python literal-tests/clause-visits/UP/branch events, NOT native FLMV TM transitions.
The complete source-indexed 660KB raw JSON, full Python producer, separate checker, full five-page PDF and checksum manifest are in the local reproducibility bundle; GitHub compact code/summary are separately verified. Native preflight and actual two-sided native FLMV timing are pending, and there is no industrial CDCL comparison.
\section{Corpus, duplication, and remaining arrow}
Surveyed current governance REDIRECT1, CONTRACT1, AFFINE-SMALLBIAS-DESCENT-GATE2, CONSTRUCTION-STRENGTHEN17, ARC-LINEAR22-SATURATION1, root main.tex, AGENTS/CODEX/skill, issue 3, LSPV-1 affine-cover261, LSPV-5 residual-cover231 and virtual-input226, LSPV-4 coupled-budget260 and index-erasure261 R1, W1/S15 fixture, primary FLMV and Cube-and-Conquer/backbone literature. The logical cover requirement and restricted raw-cofactor mass obstruction are existing peer results; the direct always-SAT cover-check reduction applies a known completeness fact to the exact ARC gate. The citation scope was found to be overbroad, not the underlying FLMV 5.7 lower bound.
\noindent\textbf{OPEN:} one uniform, all-input, hereditary, both-label-correct $A\mapsto\Phi(A)$ producing genuinely power-short ordinary SAT children and a paid syntactic/source-relative cover or exact fallback within the \emph{same} strict FLMV resource region. P versus NP remains OPEN.
\begin{thebibliography}{9}
\bibitem{FLMV} L. Fortnow, R. Lipton, D. van Melkebeek, A. Viglas. \emph{Time-Space Lower Bounds for Satisfiability}. JACM 52(6), 2005. \url{https://lance.fortnow.com/papers/files/tst.pdf}.
\bibitem{backbone} L. Hemaspaandra and D. Narvaez. \emph{The Opacity of Backbones}. AAAI 2017. DOI: 10.1609/aaai.v31i1.11134.
\bibitem{cube} M. Heule, O. Kullmann, S. Wieringa, A. Biere. \emph{Cube and Conquer}. HVC 2011. DOI: 10.1007/978-3-642-34188-5\_8.
\end{thebibliography}
\end{document}
% END_OF_VERIFIED_REPORT::LSPV-4::ARC-COVER-GATE-SCOPE262::P_VS_NP_OPEN
